\documentclass[12pt]{article}

\usepackage{tikz}
\usetikzlibrary{decorations.markings}

%\usepackage{geometry}
%\geometry{paperwidth=170mm, paperheight=240mm, left=42pt, top=40pt, textwidth=280pt, marginparsep=20pt, marginparwidth=100pt, textheight=560pt, footskip=40pt}

%\documentclass[12pt]{amsart}
\usepackage[utf8]{inputenc}
\usepackage{times}
\usepackage{geometry}
\usepackage{tabularx}
\usepackage{amsmath}
\usepackage{amsfonts}
\usepackage{amsthm}
\usepackage{amssymb}
\usepackage{stmaryrd}
\usepackage{color}
\usepackage{sidecap}
\usepackage{physics}
\usepackage{graphicx}
\usepackage{color}
\usepackage{verse}
\usepackage{hyperref}
\usepackage{tensor}
\usepackage{marginnote}
\usepackage{wrapfig}
\usepackage[font={small}]{caption}

\newcommand{\gke}{\texttt{Gkeyll}}

\newcommand{\attrib}[1]{
\nopagebreak{\raggedleft\footnotesize #1\par}\vskip0.1in}
\renewcommand{\poemtitlefont}{\normalfont\large\itshape\centering}

\newtheorem{proposition}{Proposition}
\newtheorem{theorem}{Theorem}
\newtheorem{lemma}{Lemma}
\newtheorem{remark}{Remark}
\newtheorem{definition}{Definition}
\newtheorem{principle}{Principle}
\theoremstyle{definition}
\newtheorem{exmp}{Example}

\theoremstyle{definition}
\newtheorem{entry}{Entry}

\theoremstyle{definition}
\newtheorem{exer}{Exercise}

% auto-scaled figured
\newcommand{\incfig}{\centering\includegraphics}
\setkeys{Gin}{width=0.9\linewidth,keepaspectratio}

% Commonly used macros
\newcommand{\eqr}[1]{Eq.\thinspace(#1)}
\newcommand{\pfrac}[2]{\frac{\partial #1}{\partial #2}}
\newcommand{\pfracc}[2]{\frac{\partial^2 #1}{\partial #2^2}}
\newcommand{\pfraccc}[1]{\frac{\partial^2 }{\partial #1^2}}
\newcommand{\pfraca}[1]{\frac{\partial}{\partial #1}}
\newcommand{\pfracb}[2]{\partial #1/\partial #2}
\newcommand{\pfracbb}[2]{\partial^2 #1/\partial #2^2}
\newcommand{\spfrac}[2]{{\partial_{#1}} {#2}}
\newcommand{\mvec}[1]{\mathbf{#1}}
\newcommand{\bmvec}[1]{\breve{\mathbf{#1}}}
\newcommand{\gvec}[1]{\boldsymbol{#1}}
\newcommand{\script}[1]{\mathpzc{#1}}
\newcommand{\gcn}{\nabla}
\newcommand{\gcs}{\nabla_{\mvec{x}}}
\newcommand{\gvs}{\nabla_{\mvec{v}}}

\newcommand{\resetall}{\setcounter{entry}{0}\setcounter{equation}{0}\setcounter{footnote}{0}}

\newcommand{\cbas}[1]{\gvec{\sigma}_{#1}}
\newcommand{\xbas}{\gvec{\sigma}_{1}}
\newcommand{\ybas}{\gvec{\sigma}_{2}}
\newcommand{\zbas}{\gvec{\sigma}_{3}}

\newcommand{\basis}[1]{\mvec{e}_{#1}}
\newcommand{\bbasis}[1]{\breve{\mvec{e}}_{#1}}
\newcommand{\dbasis}[1]{\mvec{e}^{#1}}
\newcommand{\bdbasis}[1]{\breve{\mvec{e}}^{#1}}

\newcommand{\nbasis}[1]{\hat{\mvec{e}}_{#1}}
\newcommand{\ndbasis}[1]{\hat{\mvec{e}}^{#1}}

\newcommand{\gbasis}[2]{\mvec{#1}_{#2}}
\newcommand{\gdbasis}[2]{\mvec{#1}^{#2}}

\newcommand{\veps}{\gvec{\varepsilon}}
\newcommand{\bdum}{\breve{\_}}
\newcommand{\vecspace}{\mathcal{V}}
\newcommand{\tvecspace}{T_pM}
\newcommand{\vnorm}[1]{\lVert{#1}\rVert}

\newcommand{\gsel}[2]{\langle{#1}\rangle_{#2}}
\newcommand{\gzsel}[1]{\langle {#1} \rangle}

\newcommand{\uln}[1]{\underline{#1}}

\newcommand{\dprod}{\mathfrak{D}}

\newcommand{\dvol}{\thinspace d^3\mvec{x}}
\newcommand{\dsurf}{\thinspace ds}

% Make the items smaller
\newcommand{\cramplist}{
	\setlength{\itemsep}{0in}
	\setlength{\partopsep}{0in}
	\setlength{\topsep}{0in}}
\newcommand{\cramp}{\setlength{\parskip}{.5\parskip}}
\newcommand{\zapspace}{\topsep=0pt\partopsep=0pt\itemsep=0pt\parskip=0pt}

%\marginpar{\raggedleft\footnotesize Margin note}

\title{Taxonomy and Construction of Base Schemes for Hyperbolic PDEs}%
\author{Ammar H. Hakim}%
\date{\today}

\begin{document}
\maketitle

\tableofcontents

\section{Exact Solution to 1D Linear Hyperbolic System}

Consider the one-dimensional system of equations
\begin{align}
\pfrac{Q}{t} + \pfrac{F}{x} = 0 \label{eq:hyp1}
\end{align}
Where $Q(x,t)$ is a vector of $m$ conserved quantities and $F(Q)$ is
the flux function. At first, let us assume the system is linear and
write the flux as
\begin{align}
F = A Q
\end{align}
where $A$ is a constant $m\times m$ matrix. We will assume that this
system is hyperbolic, i.e. the eigenvalues of $A$ are real and the
eigenvectors are complete. Let $\lambda_p$ be the eigenvalues and
$r^p$ and $l^p$, $p=1,\ldots,m$ be the right and left eigenvectors
respectively. We will represent right eigenvectors as column vectors,
and left eigenvectors as row vectors.

To solve \eqr{\ref{eq:hyp1}} we first convert it to a system of
uncoupled advection equations by multiplying it from the left by
$l^p$. This gives
\begin{align}
\pfrac{w^p}{t} + \lambda_p \pfrac{w^p}{x} = 0 \label{eq:wp-advect}
\end{align}
where we have defined the \emph{Riemann variables} $w^p = l^p Q$. Note
that given $w^p$ we can recompute $Q$ from
\begin{align}
Q = \sum_p w^p r^p. \label{eq:wp-to-Q}
\end{align}
A useful identity is
\begin{align}
\sum_p w^p \lambda_p r^p = \sum_p w^p A r^p = A \sum_p w^p r^p = AQ = F(Q). \label{eq:wp-ident1}
\end{align}

Now consider a domain $-\infty < x < \infty$ and the initial
conditions $w^p_0(x) = l^p Q_0(x)$. Each advection equation for the
Riemann variables can be solved exactly as
\begin{align}
w^p(x,t) = w^p_0(x-\lambda_p t)
\end{align}
From this, the exact solution to the linear system can be obtained as
\begin{align}
Q(x,t) = \sum_p w^p_0(x-\lambda_p t) r^p = \sum_p l^p Q_0(x-\lambda_p t) r^p.
\end{align}

\section{Basic Formulation of Finite-Volume Schemes}

The finite-volume (FV) scheme for the system \eqr{\ref{eq:hyp1}} is a
method to update the \emph{cell-average} of the solution in
time. Consider a cell $I_j \equiv [x_{j-1/2},x_{j+1/2}]$ with uniform
cell-spacing $\Delta x \equiv x_{j+1/2}-x_{j-1/2}$. Integrate
\eqr{\ref{eq:hyp1}} to get
\begin{align}
  \pfrac{Q_j}{t} + \frac{F_{j+1/2}-F_{j-1/2}}{\Delta x} = 0
  \label{eq:fv-exact}
\end{align}
where $Q_j$ are cell-average quantities
\begin{align}
  Q_j(t) = \frac{1}{\Delta x}\int_{I_j} Q(x,t) \thinspace dx
\end{align}
and $F_{j\pm 1/2} = F(Q_{j\pm 1/2})$ are the fluxes at the
\emph{cell-edges} $x_{j\pm 1/2}$. Notice that in effect finite-volume
scheme uses the \emph{mean flux gradient} in the cell
\begin{align}
  \frac{1}{\Delta x}\int_{I_j} \pfrac{F}{x} \thinspace dx =
  \frac{F_{j+1/2}-F_{j-1/2}}{\Delta x}.
\end{align}
to update the \emph{cell average} in time. This is a subtle point and
often the cause of misinterpreting the accuracy or order of a FV
scheme.

\begin{figure}
  \setkeys{Gin}{width=0.45\linewidth,keepaspectratio}
  \incfig{fv-lr-vals.pdf}
  \caption{Recovery is used to compute left/right values
    $Q_{i+1/2}^{\pm}$ from a set of cell-averages and are then fed
    into a numerical flux function to update the solution.}
  \label{fig:fv-lr}
\end{figure}

At this point the discrete expression is exact, but only formally:
given the cell-averages we can't uniquely determine the edge values
$Q_{j\pm 1/2}$ to insert into the edge fluxes. The FV scheme is an
approximation in which these edge values are \emph{recovered}
(approximately) from the cell-averages and then used in a
\emph{numerical-flux} to update the cell-average approximation. We can
write this as
\begin{align}
  \pfrac{Q_j}{t} + \frac{G(Q_{j+1/2}^+,Q_{j+1/2}^-) - G(Q_{j-1/2}^+,Q_{j-1/2}^-)}{\Delta x} = 0
\end{align}
where $G(Q_L,Q_R)$ the \emph{numerical-flux} function and
$Q_{j\pm 1/2}^+$ and $Q_{j\pm 1/2}^-$ are recovered values just the
right and left of the interface $j\pm 1/2$ respectively (see
Fig.\thinspace\ref{fig:fv-lr}). For consistency with the exact form
\eqr{\ref{eq:fv-exact}} we must ensure that the numeric flux is
Lipschitz continuous\footnote{This ensures that the derivative of the
  numerical-flux with each of its independent variables is
  \emph{bounded}.} and \emph{consistent}: when $Q_L = Q_R$ then
$G(Q_L,Q_R)$ must reduce to the physical flux function, i.e.
\begin{align}
  \lim_{Q_{L,R}\rightarrow Q} G(Q_L,Q_R) = F(Q).
\end{align}

\section{Computing the Numerical Flux}

Let us first consider a linear system in which we use the
cell-averages directly as the left/right edge values, skipping the
recovery step completely. This \emph{first-order} scheme for linear
equations will form the basis for higher-order schemes for nonlinear
systems. We will also use a simple forward-Euler time-stepper to
simplify the third-step.

Instead of directly discretizing the coupled system of equations (in
which upwind direction may not not be clear), we can instead solve the
linear advection equations for the Riemann variables using an upwind
method and the convert the solution for $w^p$ to $Q$ using
\eqr{\ref{eq:wp-to-Q}}. We can write a first-order upwind method for
the Riemann variables as
\begin{align}
w^{p,n+1}_j = w^{p,n}_j - \frac{\Delta t}{\Delta x}
\left(
G^p(w^{p,n}_{j+1},w^{p,n}_{j}) - G^p(w^{p,n}_{j},w^{p,n}_{j-1})
\right)
\end{align}
where the numerical flux function for the Riemann variables,
$G^p(w_R,w_L)$, is defined as
\begin{align}
G^p(w_R^p,w_L^p) = \frac{\lambda_p}{2}(w_R^p + w_L^p) - \frac{|\lambda_p|}{2}(w_R^p - w_L^p).
\end{align}
Note that this choice of flux ensures that if $\lambda_p>0$ then
$G^p(w_R,w_L) = \lambda_p w_L$ and if $\lambda_p < 0$ then
$G^p(w_R,w_L) = \lambda_p w_R$, ensuring proper upwinding of the
values at cell interfaces.

To convert this to a scheme for $Q$ instead, multiply by $r^p$ and sum
over $p$ to get
\begin{align}
Q^{n+1}_j = Q^n_j - \frac{\Delta t}{\Delta x}\left(G(Q^n_{j+1},Q^n_{j}) - G(Q^n_{j},Q^n_{j-1})\right).
\end{align}
The numerical flux $G(Q_R,Q_L)$ is computed from
\begin{align}
G(Q_R,Q_L) = \sum_p r^p G^p(w_R^p,w_L^p) = \frac{1}{2} \sum_p \lambda_p r^p (w_R^p+w_L^p) 
- \frac{1}{2} \sum_p |\lambda_p| r^p (w_R^p-w_L^p).
\end{align}
We can write the first term, using identity \eqr{\ref{eq:wp-ident1}}
as $(F(Q_R)+F(Q_L))/2$. To rewrite the second term introduce
\begin{align}
\lambda_p^+ &= \max(\lambda_p,0) \\
\lambda_p^- &= \min(\lambda_p,0).
\end{align}
Note that in terms of these we can write
$\lambda_p = \lambda_p^+ + \lambda_p^-$ and
$|\lambda_p| = \lambda_p^+ - \lambda_p^-$. Using the latter identity
the numerical flux can be written as
\begin{align}
G(Q_R,Q_L) = \frac{1}{2}\big(F(Q_R)+F(Q_L)\big) - \frac{1}{2}(A^+\Delta Q_{R,L} - A^-\Delta Q_{R,L})
\end{align}
where the \emph{fluctuations} $A^\pm\Delta Q$ are defined as
\begin{align}
A^\pm\Delta Q_{R,L} \equiv \sum_p r^p \lambda^\pm_p (w_R^p-w_L^p) = \sum_p r^p \lambda^\pm_p l^p(Q_R-Q_L).
\label{eq:fluct}
\end{align}
The fluctuations satisfy the \emph{flux-difference} or
\emph{flux-jump} identity
\begin{align}
A^+\Delta Q_{R,L} + A^-\Delta Q_{R,L} = 
\sum_p r^p \underbrace{(\lambda_p^+ + \lambda_p^-)}_{\lambda_p}
(w^p_R-w^p_L) = F(Q_R)-F(Q_L). \label{eq:flux-jump}
\end{align}
Using this identity, the numerical flux can also be written as
\begin{align}
G(Q_R,Q_L) = F(Q_L) + A^-\Delta Q_{R,L} = F(Q_R) - A^+\Delta Q_{R,L}.
\end{align}

For further use, we will introduce the concept of a \emph{wave}. We
define these vectors as
\begin{align}
  \mathcal{W}_{R,L}^p \equiv r^p (w_R^p-w_L^p) =  r^p  l^p(Q_R-Q_L) .
\end{align}
In terms of the waves the fluctuations, \eqr{\ref{eq:fluct}}, can be
computed as
\begin{align}
  A^\pm\Delta Q_{R,L} = \sum_p \lambda^\pm_p \mathcal{W}_{R,L}^p.
\end{align}

All the above manipulations were done for a system of linear
hyperbolic PDEs, in which the flux Jacobian does not depend on the
solution itself. However, for \emph{nonlinear} hyperbolic PDEs the
flux Jacobian \emph{does} depend on the solution. Now consider again
the two values $Q_{R,L}$ at each interface. Which one should we use to
compute the flux Jacobian and from that the eigenvalues and the
right/left eigenvectors?

\section{The $f$-wave Approach}





\end{document}

